\documentclass[a4paper,12pt]{article}
\usepackage[utf8]{inputenc}
\usepackage[T1]{fontenc}
\usepackage[ngerman]{babel}
\usepackage{amsmath, amssymb}
\usepackage{geometry}
\usepackage{tikz}
\usepackage{tikz-cd}
\usetikzlibrary{shapes.geometric, arrows.meta, positioning, fit, backgrounds}
\usepackage{parskip}
\usepackage{titlesec}
\usepackage{booktabs}
\usepackage{enumitem}

\geometry{left=3cm, right=3cm, top=2.5cm, bottom=2.5cm}

\titleformat{\section}{\large\bfseries}{}{0em}{}[\titlerule]
\titleformat{\subsection}{\normalsize\bfseries}{}{0em}{}

\tikzset{
  box/.style={rectangle, draw, rounded corners=3pt, text width=4.5cm,
              align=center, minimum height=1.0cm, font=\small, inner sep=4pt},
  boxw/.style={box, text width=5cm},
  arrow/.style={-{Latex[length=2mm]}, thick},
}

\begin{document}

\begin{center}
  {\LARGE\bfseries Studierhinweise zu Lektion 1}\\[0.5em]
  {\large Mathematische Grundlagen (Modul 61111)}
\end{center}

\vspace{1em}

Diese Lektion legt die sprachlichen und begrifflichen Grundlagen für den gesamten
weiteren Verlauf des Moduls. Im ersten Kapitel werden Notation und Grundbegriffe
eingeführt, die in der Mathematik allgegenwärtig sind: das Summensymbol~$\Sigma$,
logische Junktoren und Quantoren, Mengen und Abbildungen, Verknüpfungen und
Körper. In Kapitel~2 wird die Matrizenrechnung eingeführt, in Kapitel~3 werden
Elementarmatrizen und elementare Zeilenumformungen behandelt.

Einen mathematischen Text darf man nicht wie einen Roman lesen -- man muss ihn
sich \emph{erarbeiten}. Schreiben Sie mit, stellen Sie sich Beispiele vor, und lösen
Sie die Übungsaufgaben aktiv, bevor Sie in die Lösungen schauen.

% -----------------------------------------------------------------------
\section*{Struktur der Lektion 1}

\begin{center}
\begin{tikzpicture}[node distance=0.6cm and 0.8cm]

  % Column headers
  \node[font=\bfseries\small] (h1) at (0,0)    {Kapitel 1: Grundbegriffe};
  \node[font=\bfseries\small] (h2) at (6.5,0)  {Kapitel 2: Matrizen};
  \node[font=\bfseries\small] (h3) at (11.5,0) {Kapitel 3: Elementarmatrizen};

  % Chapter 1 boxes
  \node[box, below=0.8cm of h1] (sum)  {1.1 Summensymbol~$\Sigma$\\Doppelsummen};
  \node[box, below=0.5cm of sum] (aus) {1.2 Aussagen\\Junktoren, Quantoren,\\Vollst. Induktion};
  \node[box, below=0.5cm of aus] (men) {1.3 Mengen\\Vereinigung, Schnitt,\\Produkt, Mächtigkeit};
  \node[box, below=0.5cm of men] (abb) {1.4 Abbildungen\\Bild, Urbild,\\injektiv, surjektiv, bijektiv};
  \node[box, below=0.5cm of abb] (vkn) {1.5 Verknüpfungen\\kommutativ, assoziativ,\\neutrales Element};
  \node[box, below=0.5cm of vkn] (koe) {1.6 Körper\\$K$ mit $+$, $\cdot$,~$0$,~$1$};

  % Chapter 2 boxes
  \node[box, below=0.8cm of h2] (mat) {2.1--2.2 Matrizen\\Addition, Skalarmultiplikation};
  \node[box, below=0.5cm of mat] (mul) {2.3 Matrizenmultiplikation\\Einheitsmatrix $I_m$};
  \node[box, below=0.5cm of mul] (gem) {2.4 Gemischte Rechenregeln\\Distributivgesetze};

  % Chapter 3 boxes
  \node[box, below=0.8cm of h3] (elz) {3.1 Elementare Zeilenumformungen\\$Z_{ij}$, $Z_i(r)$, $Z_{ij}(s)$};
  \node[box, below=0.5cm of elz] (elm) {3.2 Elementarmatrizen\\$P_{ij}$, $D_i(r)$, $T_{ij}(s)$};
  \node[box, below=0.5cm of elm] (zeq) {3.3 Zeilenäquivalenz\\$A \sim_Z B$};

  % Arrows within Kap 1
  \draw[arrow] (sum) -- (aus);
  \draw[arrow] (aus) -- (men);
  \draw[arrow] (men) -- (abb);
  \draw[arrow] (abb) -- (vkn);
  \draw[arrow] (vkn) -- (koe);

  % Arrows within Kap 2
  \draw[arrow] (mat) -- (mul);
  \draw[arrow] (mul) -- (gem);

  % Arrows within Kap 3
  \draw[arrow] (elz) -- (elm);
  \draw[arrow] (elm) -- (zeq);

  % Cross-chapter arrows
  \draw[arrow] (koe.east) -- ++(0.3,0) |- (mat.west);
  \draw[arrow] (mul.east) -- ++(0.3,0) |- (elz.west);

\end{tikzpicture}
\end{center}

\newpage

% -----------------------------------------------------------------------
\section*{Zielelement 1.1 -- Das Summensymbol $\Sigma$}

\subsection*{Lerninhalte}

\begin{center}
\begin{tikzpicture}[node distance=0.6cm and 1.0cm]
  \node[box] (def)  {Definition des Summensymbols $\displaystyle\sum_{i=a}^{b} f(i)$};
  \node[box, below=0.5cm of def] (bsp) {Verschiedene Schreibweisen und Indexvariablen};
  \node[box, below=0.5cm of bsp] (dop) {Doppelsummen $\displaystyle\sum_{i}\sum_{j} a_{ij}$};
  \node[box, below=0.5cm of dop] (mrk) {Merkregel: Vertauschen von Summensymbolen};
  \draw[arrow] (def) -- (bsp);
  \draw[arrow] (bsp) -- (dop);
  \draw[arrow] (dop) -- (mrk);
\end{tikzpicture}
\end{center}

\subsection*{Lernziele}

Nach Durcharbeiten dieses Abschnitts sollten Sie:
\begin{itemize}
  \item das Summensymbol $\Sigma$ lesen, schreiben und ausrechnen können,
  \item Summen in Sigma-Schreibweise umwandeln und umgekehrt,
  \item keine Angst vor Doppelsummen haben und die Merkregel 1.1.1 anwenden
        können (Vertauschen von Summationsreihenfolgen).
\end{itemize}

\subsection*{Selbstkontrollelement 1.1}

Berechnen Sie $\displaystyle\sum_{i=1}^{4}\sum_{j=1}^{3} (i+j)$, einmal indem Sie
zuerst über $j$ und dann über $i$ summieren, und einmal in umgekehrter Reihenfolge.
Erhalten Sie dasselbe Ergebnis?

\newpage

% -----------------------------------------------------------------------
\section*{Zielelement 1.2 -- Aussagen, Junktoren und Quantoren}

\subsection*{Lerninhalte}

\begin{center}
\begin{tikzpicture}[node distance=0.6cm and 1.0cm]
  \node[box] (aus) {Aussage: wahr oder falsch};
  \node[box, right=1.0cm of aus] (jun) {Junktoren:\\$\lnot,\,\land,\,\lor,\,\Rightarrow,\,\Leftrightarrow$};
  \node[box, below=0.5cm of aus] (wt)  {Wahrheitstafeln};
  \node[box, below=0.5cm of wt]  (leq) {Logische Äquivalenz};
  \node[box, below=0.5cm of leq] (qua) {Quantoren: $\exists$, $\forall$};
  \node[box, below=0.5cm of qua] (neg) {Negation von\\All- und Existenzaussagen};
  \node[box, below=0.5cm of neg] (ind) {Vollständige Induktion\\(Induktionsanfang,\\Induktionsschritt)};
  \node[box, below=0.5cm of ind] (bew) {Beweisprinzipien:\\direkt, Kontraposition,\\Widerspruch, Ringschluss};
  \draw[arrow] (aus) -- (wt);
  \draw[arrow] (jun.south) |- (wt.east);
  \draw[arrow] (wt) -- (leq);
  \draw[arrow] (leq) -- (qua);
  \draw[arrow] (qua) -- (neg);
  \draw[arrow] (neg) -- (ind);
  \draw[arrow] (ind) -- (bew);
\end{tikzpicture}
\end{center}

\subsection*{Lernziele}

Nach Durcharbeiten dieses Abschnitts sollten Sie:
\begin{itemize}
  \item die Wahrheitstafeln der fünf Junktoren kennen und aufstellen können,
  \item nicht übermäßig komplexe Aussagen mit Quantoren $\exists$, $\forall$ ausdrücken
        und in Umgangssprache übersetzen können,
  \item die Negation von All- und Existenzaussagen bilden können (inkl.\ verschachtelter Quantoren),
  \item das Prinzip der vollständigen Induktion verstehen und in Beweisen anwenden können,
  \item die verschiedenen Beweisprinzipien (direkt, Kontraposition, Widerspruch, Ringschluss) kennen und erkennen.
\end{itemize}

\subsection*{Selbstkontrollelement 1.2}

Bestimmen Sie die Negation der Aussage:
\[
  \forall\,\varepsilon > 0 \; \exists\,n_\varepsilon \in \mathbb{N} \;
  \forall\,n \geq n_\varepsilon : |a_n - a| < \varepsilon.
\]

\newpage

% -----------------------------------------------------------------------
\section*{Zielelement 1.3 -- Mengen}

\subsection*{Lerninhalte}

\begin{center}
\begin{tikzpicture}[node distance=0.6cm and 1.0cm]
  \node[box] (mng) {Menge, Element ($\in$, $\notin$)};
  \node[box, below=0.5cm of mng] (tlm) {Teilmenge ($\subseteq$), Gleichheit von Mengen};
  \node[box, below=0.5cm of tlm] (ops) {Mengenoperationen:\\$\cup$,\; $\cap$,\; $\setminus$,\; $\emptyset$,\; disjunkt};
  \node[box, below=0.5cm of ops] (end) {Endliche / unendliche Mengen\\Mächtigkeit $|M|$};
  \node[box, below=0.5cm of end] (pro) {Produktmenge $M \times N$\\geordnetes Paar $(m,n)$};
  \draw[arrow] (mng) -- (tlm);
  \draw[arrow] (tlm) -- (ops);
  \draw[arrow] (ops) -- (end);
  \draw[arrow] (end) -- (pro);
\end{tikzpicture}
\end{center}

\subsection*{Lernziele}

Nach Durcharbeiten dieses Abschnitts sollten Sie:
\begin{itemize}
  \item mit den Begriffen Vereinigung, Durchschnitt, Differenzmenge, Gleichheit
        und Mächtigkeit von Mengen sicher umgehen können,
  \item die Produktmenge $M \times N$ und geordnete Paare erklären können,
  \item einfache Beweise über Mengen (z.\,B.\ zur Mächtigkeit von $M \cup N$)
        selbstständig führen können.
\end{itemize}

\subsection*{Selbstkontrollelement 1.3}

Seien $M = \{1,2,3,4\}$ und $N = \{3,4,5,6\}$.
Bestimmen Sie $M \cup N$, $M \cap N$, $M \setminus N$ und $|M \times N|$.

\newpage

% -----------------------------------------------------------------------
\section*{Zielelement 1.4 -- Abbildungen}

\subsection*{Lerninhalte}

\begin{center}
\begin{tikzpicture}[node distance=0.6cm and 1.0cm]
  \node[box] (def) {Abbildung $f: M \to N$\\als Teilmenge von $M \times N$};
  \node[box, below=0.5cm of def] (bil) {Bild $f(m)$, Bildbereich $\mathrm{Bild}(f)$};
  \node[box, below=0.5cm of bil] (urb) {Urbild};
  \node[box, below=0.5cm of urb] (inj) {injektiv, surjektiv, bijektiv};
  \node[box, below=0.5cm of inj] (idb) {Identische Abbildung $\mathrm{id}_M$};
  \node[box, below=0.5cm of idb] (kom) {Komposition $g \circ f$\\(Assoziativgesetz)};
  \node[box, below=0.5cm of kom] (inv) {Invertierbarkeit\\$\Leftrightarrow$ bijektiv (Kor. 1.4.22)};
  \draw[arrow] (def) -- (bil);
  \draw[arrow] (bil) -- (urb);
  \draw[arrow] (urb) -- (inj);
  \draw[arrow] (inj) -- (idb);
  \draw[arrow] (idb) -- (kom);
  \draw[arrow] (kom) -- (inv);
\end{tikzpicture}
\end{center}

\subsection*{Lernziele}

Nach Durcharbeiten dieses Abschnitts sollten Sie:
\begin{itemize}
  \item den Abbildungsbegriff präzise formulieren und von der schulischen
        Beschreibung abgrenzen können,
  \item Bild und Urbilder von Abbildungen bestimmen können,
  \item die Begriffe injektiv, surjektiv und bijektiv erklären und überprüfen können,
  \item die Komposition von Abbildungen ausführen und das Assoziativgesetz anwenden,
  \item Korollar~1.4.22 (Charakterisierung invertierbarer Abbildungen) kennen und
        einsetzen können.
\end{itemize}

\subsection*{Selbstkontrollelement 1.4}

Sei $f: \mathbb{Z} \to \mathbb{Z}$ definiert durch $f(n) = 2n$.
Ist $f$ injektiv? Ist $f$ surjektiv? Begründen Sie Ihre Antwort.

\newpage

% -----------------------------------------------------------------------
\section*{Zielelement 1.5 -- Verknüpfungen}

\subsection*{Lerninhalte}

\begin{center}
\begin{tikzpicture}[node distance=0.6cm and 1.0cm]
  \node[box] (vkn) {Verknüpfung auf $M$:\\Abbildung $M \times M \to M$};
  \node[box, below=0.5cm of vkn] (eig) {Eigenschaften:\\kommutativ, assoziativ,\\neutrales Element};
  \node[box, below=0.5cm of eig] (inv) {Invertierbarkeit\\bezüglich der Verknüpfung};
  \node[box, below=0.5cm of inv] (dis) {Distributivgesetze\\(zwei Verknüpfungen)};
  \node[box, below=0.5cm of dis] (f2)  {Beispiel: $\mathbb{F}_2 = \{0,1\}$\\mit $+$ und $\cdot$};
  \draw[arrow] (vkn) -- (eig);
  \draw[arrow] (eig) -- (inv);
  \draw[arrow] (inv) -- (dis);
  \draw[arrow] (dis) -- (f2);
\end{tikzpicture}
\end{center}

\subsection*{Lernziele}

Nach Durcharbeiten dieses Abschnitts sollten Sie:
\begin{itemize}
  \item den Begriff der Verknüpfung auf einer Menge als Abbildung verstehen,
  \item die Eigenschaften kommutativ, assoziativ, neutrales Element und
        Invertierbarkeit an konkreten Beispielen überprüfen können,
  \item den Körper $\mathbb{F}_2$ in- und auswendig kennen.
\end{itemize}

\subsection*{Selbstkontrollelement 1.5}

Ist die Komposition $\circ$ auf $\mathrm{Abb}(M,M)$ kommutativ?
Geben Sie ein Gegenbeispiel oder einen Beweis an.

\newpage

% -----------------------------------------------------------------------
\section*{Zielelement 1.6 -- Körper}

\subsection*{Lerninhalte}

\begin{center}
\begin{tikzpicture}[node distance=0.6cm and 1.0cm]
  \node[box] (def) {Körper $K$: Menge mit $+$, $\cdot$, $0$, $1$};
  \node[box, below=0.5cm of def] (axm) {Axiome: Kommutativität, Assoziativität,\\neutrales Element, Invertierbarkeit,\\Distributivgesetz};
  \node[box, below=0.5cm of axm] (bsp) {Beispiele: $\mathbb{Q}$, $\mathbb{R}$, $\mathbb{C}$, $\mathbb{F}_2$};
  \node[box, below=0.5cm of bsp] (rec) {Rechenregeln (Prop.\ 1.6.4):\\$a \cdot 0 = 0$, Nullteilerfreiheit,\\Eindeutigkeit inverser Elemente};
  \node[box, below=0.5cm of rec] (not) {Notation: $-a$, $a^{-1}$, $\frac{b}{a}$};
  \draw[arrow] (def) -- (axm);
  \draw[arrow] (axm) -- (bsp);
  \draw[arrow] (bsp) -- (rec);
  \draw[arrow] (rec) -- (not);
\end{tikzpicture}
\end{center}

\subsection*{Lernziele}

Nach Durcharbeiten dieses Abschnitts sollten Sie:
\begin{itemize}
  \item die Definition eines Körpers auswendig kennen und an Beispielen überprüfen können,
  \item verstehen, warum $\mathbb{Z}$ kein Körper ist,
  \item die Rechenregeln in Proposition~1.6.4 aus den Körperaxiomen herleiten können,
  \item mit $\mathbb{F}_2$ sicher rechnen können.
\end{itemize}

\subsection*{Selbstkontrollelement 1.6}

Beweisen Sie: In einem Körper $K$ gilt $(-1) \cdot (-1) = 1$.
\emph{Hinweis:} Verwenden Sie Bemerkung~1.6.6.

\newpage

% -----------------------------------------------------------------------
\section*{Zielelement 2 -- Matrizen}

\subsection*{Lerninhalte}

\begin{center}
\begin{tikzpicture}[node distance=0.6cm and 1.0cm]
  \node[box] (def) {$m \times n$-Matrix über $K$\\Eintrag $a_{ij}$, Zeile, Spalte};
  \node[box, below=0.5cm of def] (add) {Matrizenaddition,\\Nullmatrix, Skalarmultiplikation};
  \node[box, below=0.5cm of add] (mul) {Matrizenmultiplikation $(AB)_{ik} = \sum_j a_{ij}b_{jk}$};
  \node[box, below=0.5cm of mul] (ein) {Einheitsmatrix $I_m$,\\Assoziativgesetz};
  \node[box, below=0.5cm of ein] (inv) {Invertierbare Matrix $A^{-1}$,\\$A^{-1} A = A A^{-1} = I_n$};
  \draw[arrow] (def) -- (add);
  \draw[arrow] (add) -- (mul);
  \draw[arrow] (mul) -- (ein);
  \draw[arrow] (ein) -- (inv);
\end{tikzpicture}
\end{center}

\subsection*{Lernziele}

Nach Durcharbeiten dieses Kapitels sollten Sie:
\begin{itemize}
  \item Matrizen addieren und (wenn die Dimensionen passen) multiplizieren können,
  \item wissen, wann ein Produkt $AB$ definiert ist, und die Dimensionen des
        Ergebnisses bestimmen können,
  \item die Matrizenmultiplikation als nicht kommutative Verknüpfung auf $M_{nn}(K)$
        einordnen können,
  \item den Begriff der invertierbaren Matrix kennen.
\end{itemize}

\subsection*{Selbstkontrollelement 2}

Seien $A = \begin{pmatrix}1 & 2 \\ 0 & 1\end{pmatrix}$ und
$B = \begin{pmatrix}1 & 0 \\ 3 & 1\end{pmatrix}$.
Berechnen Sie $AB$ und $BA$. Sind die Ergebnisse gleich?

\newpage

% -----------------------------------------------------------------------
\section*{Zielelement 3 -- Elementarmatrizen}

\subsection*{Lerninhalte}

\begin{center}
\begin{tikzpicture}[node distance=0.6cm and 1.0cm]
  \node[box] (elz) {Elementare Zeilenumformungen:\\$Z_{ij}$: Zeilen tauschen\\$Z_i(r)$: Zeile skalieren\\$Z_{ij}(s)$: Zeile addieren};
  \node[box, below=0.5cm of elz] (elm) {Elementarmatrizen $P_{ij}$, $D_i(r)$, $T_{ij}(s)$\\als Multiplikation von links};
  \node[box, below=0.5cm of elm] (inv) {Elementarmatrizen sind invertierbar,\\Inverse sind wieder Elementarmatrizen};
  \node[box, below=0.5cm of inv] (zeq) {Zeilenäquivalenz $A \sim_Z B$};
  \draw[arrow] (elz) -- (elm);
  \draw[arrow] (elm) -- (inv);
  \draw[arrow] (inv) -- (zeq);
\end{tikzpicture}
\end{center}

\subsection*{Lernziele}

Nach Durcharbeiten dieses Kapitels sollten Sie:
\begin{itemize}
  \item die drei Typen elementarer Zeilenumformungen kennen und anwenden können,
  \item wissen, welche Elementarmatrix welche elementare Zeilenumformung bewirkt
        (Merkregel~3.2.8),
  \item verstehen, dass Elementarmatrizen invertierbar sind und ihre Inversen
        ebenfalls Elementarmatrizen sind,
  \item den Begriff der Zeilenäquivalenz erklären können.
\end{itemize}

\subsection*{Selbstkontrollelement 3}

Geben Sie die Elementarmatrizen $P_{12}$, $D_2(3)$ und $T_{21}(-2)$ in $M_{33}(\mathbb{R})$
explizit an. Welches sind ihre Inversen?

\end{document}
